\section{توافق جذر العدد والمشتقة اللوغاريتمية}

تبدأ الدفعة الثالثة بإغلاق الحقائق التي تسبق برهان عدم الانعدام عند
\(s=1\). هذه الحقائق ضرورية، لكنها لا تكفي وحدها لإثبات
\(L(1,\chi)\ne0\).

\begin{proposition}[توافق جذري العدد]
\label{prop:dirichlet-root-number-compatibility}
\resultid{ANT-PROP-07-05}
\provedhere
لتكن \(\chi\bmod q\) شخصية بدائية، و\(q>1\)، ولتكن
\(a\in\{0,1\}\) محددة بالعلاقة \(\chi(-1)=(-1)^a\). إذا
\[
\varepsilon_\chi=\frac{\tau(\chi)}{i^a\sqrt q},
\]
فإن
\[
\varepsilon_\chi\varepsilon_{\overline\chi}=1,
\qquad
\varepsilon_{\overline\chi}=\overline{\varepsilon_\chi}.
\]
\end{proposition}

\begin{proof}
للشخصيتين \(\chi\) و\(\overline\chi\) الزوجية نفسها \(a\). ومن مبرهنة
مجموع غاوس:
\[
\tau(\chi)\tau(\overline\chi)=\chi(-1)q=(-1)^a q.
\]
لذلك
\[
\varepsilon_\chi\varepsilon_{\overline\chi}
=
\frac{(-1)^a q}{i^{2a}q}=1,
\]
لأن \(i^{2a}=(-1)^a\). وبما أن كلا الجذرين معياره \(1\)، فإن
\(\varepsilon_{\overline\chi}=\varepsilon_\chi^{-1}\) يساوي
\(\overline{\varepsilon_\chi}\).
\end{proof}

\begin{corollary}[اتساق تطبيق المعادلة الوظيفية مرتين]
\label{cor:functional-equation-involution}
\resultid{ANT-COR-07-06}
\provedhere
تطبيق المعادلة الوظيفية مرتين يعيد الدالة المكتملة نفسها:
\[
\Lambda(s,\chi)
=
\varepsilon_\chi\varepsilon_{\overline\chi}\Lambda(s,\chi)
=
\Lambda(s,\chi).
\]
\end{corollary}

\begin{proof}
نطبق
\[
\Lambda(s,\chi)=\varepsilon_\chi\Lambda(1-s,\overline\chi)
\]
ثم نطبق المعادلة الوظيفية على الطرف الثاني، ونستعمل القضية السابقة.
\end{proof}

\begin{proposition}[المشتقة اللوغاريتمية]
\label{prop:dirichlet-l-logarithmic-derivative}
\resultid{ANT-PROP-07-06}
\provedhere
لكل شخصية ديريشليه \(\chi\bmod q\)، وعندما \(\Re(s)>1\):
\[
-\frac{L'}{L}(s,\chi)
=
\sum_{n=1}^{\infty}\frac{\Lambda(n)\chi(n)}{n^s}.
\]
وتتقارب السلسلة مطلقًا ومحليًا بانتظام في هذا النصف من المستوى.
\end{proposition}

\begin{proof}
من المنتج الأويلري وعدم الانعدام في \(\Re(s)>1\):
\[
\log L(s,\chi)
=
\sum_p\sum_{k=1}^{\infty}
\frac{\chi(p)^k}{k p^{ks}},
\]
حيث تتقارب السلسلة مطلقًا ومحليًا بانتظام. بالتفاضل حدًا حدًا:
\[
-\frac{L'}{L}(s,\chi)
=
\sum_p\sum_{k=1}^{\infty}
\frac{\chi(p)^k\log p}{p^{ks}}.
\]
وبما أن \(\chi(p)^k=\chi(p^k)\)، وأن
\(\Lambda(n)=\log p\) عندما \(n=p^k\) وتساوي صفرًا خلاف ذلك، نحصل على
الصيغة المعلنة. أما التقارب المطلق فينتج من
\[
\sum_{n=1}^{\infty}\frac{\Lambda(n)|\chi(n)|}{n^\sigma}
\leq
\sum_{n=1}^{\infty}\frac{\Lambda(n)}{n^\sigma}
<\infty
\qquad(\sigma>1).
\]
\end{proof}

\section{بوابة عدم الانعدام عند \(s=1\)}

\begin{theorem}[عدم الانعدام عند الواحد]
\label{thm:dirichlet-l-nonvanishing-at-one}
\resultid{ANT-THM-07-09}
\deferredresult{الدفعة البرهانية الثالثة من الفصل السابع}
إذا كانت \(\chi\) شخصية ديريشليه غير رئيسية، فإن
\[
L(1,\chi)\ne0.
\]
\end{theorem}

هذه النتيجة ليست مستنتجة من الهولومورفية عند \(1\). ما ثبت حتى الآن هو
وجود قيمة منتهية فقط. كما أن المنتج الأويلري يثبت عدم الانعدام في النصف
المفتوح \(\Re(s)>1\)، ولا يسمح بالمرور إلى الحد من دون حجة إضافية.

سيقسم البرهان الكامل إلى حالتين:

\begin{enumerate}
  \item \textbf{الشخصيات غير الحقيقية:} بناء دالة مساعدة أو حاصل ضرب من
  دوال \(L\) يضبط رتبة الصفر والقطب عند \(1\)، مع استعمال معاملات
  لوغاريتمية غير سالبة.
  \item \textbf{الشخصيات الحقيقية:} معالجة مستقلة تمنع الدوران المنطقي مع
  مبرهنة ديريشليه للأوليات في المتتاليات الحسابية.
\end{enumerate}

وثقت شروط القبول والمراجع في
\texttt{research/literature-reviews/chapter-07-batch-03-nonvanishing-evidence.md}.
ولا تنتقل النتيجة إلى \textenglish{PROVED-HERE} قبل إغلاق الحالتين.

\section{فرضية ريمان المعممة لدوال ديريشليه}

بعد فصل الأصفار البديهية، تقع الأصفار الأخرى في الشريط
\(0\leq\Re(s)\leq1\). وتسمى هذه الأصفار غير بديهية.

\begin{definition}[فرضية ريمان المعممة لدوال ديريشليه]
تنص فرضية ريمان المعممة في هذه العائلة على أن كل صفر غير بديهي لكل دالة
ديريشليه \(L(s,\chi)\) يحقق
\[
\Re(\rho)=\frac12.
\]
\end{definition}

هذه فرضية مفتوحة. المعادلة الوظيفية وتناظر الأصفار لا يثبتانها؛ إنهما
يثبتان فقط تناظر مجموعة الأصفار حول الخط الحرج. ولا يقدم هذا الفصل أي
ادعاء بتقدم نحو إثباتها.

\section{حالة الدفعة الثالثة}

أُغلقت في هذه المرحلة:

\begin{itemize}
  \item هوية جذري العدد واتساق المعادلة الوظيفية عند تطبيقها مرتين.
  \item المشتقة اللوغاريتمية ومتسلسلة فون مانغولت الملتوية في
  \(\Re(s)>1\).
  \item الصياغة الدقيقة لبوابة عدم الانعدام وحدودها المنطقية.
  \item صياغة فرضية ريمان المعممة بوصفها مسألة مفتوحة.
\end{itemize}

ويبقى العمل التالي هو البرهان الفعلي لعدم انعدام \(L(1,\chi)\)، لا مجرد
إعادة صياغة النتيجة أو الاستشهاد بها ضمنيًا.
